% !PRINTED 11
\weberpage{11}

%<<CONT
ergibt sich in den drei Fällen:

\[
\begin{array}{ll}
1) & \dfrac{dx}{\sqrt{-x(1-x)(1-\lambda^2x)}}, \\[10pt]
2) & \dfrac{\lambda\,dx}{\sqrt{x(1-x)(1-\lambda^2x)}}, \\[10pt]
3) & \dfrac{dx}{\sqrt{x(1-x)(1-\lambda^2x)}}.
\end{array}
\]

\begin{center}
\textbf{\S\,4. Die \gesperrt{Legendresche} Normalform.}
\end{center}

Wenn die Funktion $f(x)$, die in dem elliptischen Differential unter dem Wurzelzeichen steht, reelle Koeffizienten hat, so sind drei Fälle möglich:

1. $f(x)$ hat vier reelle Wurzeln $x_1$, $x_2$, $x_3$, $x_4$;

2. $f(x)$ hat zwei reelle Wurzeln $x_1$, $x_2$ und ein paar konjugiert imaginäre Wurzeln $x_3$, $x_4$;

3. $f(x)$ hat zwei Paare konjugiert imaginärer Wurzeln $x_1$, $x_2$ und $x_3$, $x_4$.

In allen drei Fällen läßt sich das elliptische Differential durch eine reelle lineare Transformation auf die Form bringen:

\weq{1}{\dfrac{dz}{\sqrt{(z^2-\alpha)(z^2-\beta)}},}

worin $\alpha$ und $\beta$ reelle (positive oder negative) Konstanten sind.

Wir bestimmen die lineare Abhängigkeit zwischen den Variablen $x$ und $z$ so, daß sich folgende Werte entsprechen:

\weq{2}{\begin{array}{ccccc}
x, & x_1, & x_2, & x_3, & x_4, \\
z, & \sqrt{\alpha}, & -\sqrt{\alpha}, & \sqrt{\beta}, & -\sqrt{\beta},
\end{array}}

und erhalten eine Substitution der Form:

\weq{3}{h\,\dfrac{x-x_1}{x-x_2} = \dfrac{z-\sqrt{\alpha}}{z+\sqrt{\alpha}}.}

Um $h$ zu bestimmen, setzen wir $x = x_3$, $x = x_4$ und entsprechend $z = \sqrt{\beta}$, $z = -\sqrt{\beta}$, und erhalten

\[
h\,\dfrac{(x_3x_1)}{(x_3x_2)} = \dfrac{\sqrt{\beta}-\sqrt{\alpha}}{\sqrt{\beta}+\sqrt{\alpha}}, \quad h\,\dfrac{(x_4x_1)}{(x_4x_2)} = \dfrac{\sqrt{\beta}+\sqrt{\alpha}}{\sqrt{\beta}-\sqrt{\alpha}}.
\]

Daraus durch Multiplikation und Division

\weq{4}{h = \sqrt{\dfrac{(x_3x_2)(x_4x_2)}{(x_3x_1)(x_4x_1)}}, \quad \dfrac{\sqrt{\alpha}-\sqrt{\beta}}{\sqrt{\alpha}+\sqrt{\beta}} = \sqrt{\dfrac{(x_3x_1)(x_4x_2)}{(x_3x_2)(x_4x_1)}}.}
